% !TeX program = xelatex
% ============================================================
% reflexionsformular.tex
% Eigenständiges Dokument (nicht Teil von main.tex): Handbuch Kap. 7.4.2
% (S. 92) verlangt bei KI-Einsatz "zusätzlich zur Selbständigkeitserklärung
% ein Reflexionsformular", das "bestätigt, dass die [...] Richtlinien
% eingehalten wurden, und [...] Raum [gibt], den Einsatz von KI beim
% eigenen Schreiben zu deklarieren und zu reflektieren."
%
% ACHTUNG: Das Handbuch beschreibt nur Zweck und Pflichtinhalt dieses
% Formulars, liefert aber -- anders als bei der Selbständigkeitserklärung
% (dort mit Wortlaut und Verweis auf eine Sharepoint-Vorlage) -- keine
% konkrete Layout-Vorlage im Handbuch selbst. Dies ist daher eine
% eigene, aus den Handbuch-Vorgaben abgeleitete Vorlage. Vor der Abgabe
% bei der Betreuungsperson bzw. im Sharepoint (Office 365, Info-Portal/
% Gymnasium/Maturaarbeit) prüfen, ob eine offizielle Formularvorlage
% existiert, und diese Datei bei Bedarf entsprechend anpassen.
%
% Jede/r Autor/in füllt das Formular einzeln aus (analog zur
% Selbständigkeitserklärung, die ebenfalls "ich" statt "wir" verwendet).
% Kompilieren: xelatex reflexionsformular.tex
% ============================================================

\documentclass[11pt,a4paper]{article}

\usepackage[
  a4paper,
  top=3cm, bottom=2.8cm, left=3cm, right=2cm,
]{geometry}
\usepackage{fontspec}
\setmainfont{Carlito}
\usepackage[ngerman]{babel}
\usepackage{csquotes}
\setquotestyle{swiss}
\usepackage{setspace}
\onehalfspacing
\usepackage{enumitem}
\usepackage[hidelinks]{hyperref}

\pagestyle{empty}

\begin{document}

\begin{center}
  {\LARGE\bfseries Reflexionsformular zum Einsatz von Künstlicher
  Intelligenz\par}
  \vspace{0.3em}
  {\large gemäss Handbuch Projekte KSWO 2025, Kap.\,7.4.2\par}
\end{center}

\vspace{1.5em}

\noindent
\begin{tabular}{@{}ll}
  \textbf{Titel der Arbeit:} & Bier selber brauen \\[0.3em]
  \textbf{Name der/des Verfassenden:} & \hrulefill \\[0.3em]
  \textbf{Klasse/Abteilung:} & \hrulefill \\[0.3em]
  \textbf{Datum:} & \hrulefill \\
\end{tabular}

\vspace{1.5em}

\section*{1. Eingesetzte KI-Tools}

Alle verwendeten Tools mit Datum und (zusammengefasster) Aufforderung
auflisten. Ausführliche Angaben stehen auch im Abschnitt
\enquote{Verwendete KI-Tools} im Quellenverzeichnis der Arbeit
(\texttt{ki-tools.tex}); hier reicht eine kurze Übersicht.

\begin{center}
\begin{tabular}{p{3.2cm}p{2.5cm}p{7.5cm}}
  \textbf{Tool} & \textbf{Datum} & \textbf{Wofür eingesetzt} \\
  \hline
  & & \\[2.5em]
  & & \\[2.5em]
  & & \\[2.5em]
\end{tabular}
\end{center}

\section*{2. Bestätigung}

\noindent Ich bestätige, dass
\begin{itemize}[leftmargin=1.2em]
  \item ich KI-generierte Inhalte nicht unverändert übernommen, sondern
    geprüft, bearbeitet und in eigenen Worten eingearbeitet habe,
  \item ich für die sprachliche und wissenschaftliche Qualität des
    gesamten Textes selbst verantwortlich bleibe,
  \item ich den Einsatz der oben genannten Tools vollständig und korrekt
    im Quellenverzeichnis der Arbeit deklariert habe,
  \item ich keine KI-generierten Inhalte als eigene Primär- oder
    Sekundärliteratur ausgegeben habe.
\end{itemize}

\section*{3. Reflexion}

\noindent Kurz beantworten (je 3--5 Sätze):

\begin{enumerate}[leftmargin=1.4em]
  \item Wofür genau habe ich KI eingesetzt (z.\,B. Rückmeldungen zu
    bereits Geschriebenem, Strukturierungshilfe, Sprachprüfung,
    technische Unterstützung beim LaTeX-Dokument)?

    \vspace{4em}

  \item Wie hat der Einsatz von KI meinen eigenen Schreib- bzw.
    Arbeitsprozess beeinflusst?

    \vspace{4em}

  \item Wo lagen Grenzen oder Risiken (z.\,B. fehlende/falsche
    Quellenbelege, ungenaue oder erfundene Informationen, Texte ohne
    nachvollziehbare Herkunft)?

    \vspace{4em}

  \item Was würde ich beim nächsten Mal anders machen?

    \vspace{4em}
\end{enumerate}

\vspace{2em}

\noindent
\begin{minipage}{0.45\textwidth}
  Ort, Datum \\[2.5cm]
  \hrulefill \\
  Unterschrift
\end{minipage}

\end{document}
